\chapter{Summary and outlook}
\label{ch:conclusion}
In this thesis, we present a theoretical investigation of excitonic transitions and wave-packet formation in two-dimensional transition-metal dichalcogenides under excitation by structured light, including scalar and vector vortex beams.

First, we investigate the angular-momentum structure of Laguerre--Gaussian (LG) beams and vector vortex beams (VVBs) in the angular spectrum representation. By transforming the vector potential into the Coulomb gauge, the spin--orbit interaction (SOI) naturally emerges in the longitudinal component of LG beams. For VVBs formed by the superposition of two LG beams, the resulting longitudinal angular geometry is governed by the total angular momentum difference $\Delta J=(\sigma'+\ell')-(\sigma+\ell)$, while the transverse polarization geometry is determined by the OAM difference. We demonstrate that these angular geometries persist beyond the equal-OAM condition and remain robust as long as sufficient spatial overlap exists between the constituent beams.

Next, we investigate the interaction between VVBs and valley-mixed excitons using Fermi's golden rule. The exciton transition-rate distribution in center-of-mass momentum space is determined by the compatibility between the angular-momentum differences of the optical field and the winding number of the excitonic transition dipole moments. In particular, gray excitons can be interpreted as zero-winding excitonic states due to its transition dipole moment is almost isotropic, naturally coupling to longitudinal optical structures with zero TAM difference. Matching between the optical and excitonic angular structures enables full-azimuthal momentum-selective excitation, whereas mismatching leads to directional excitation.

Finally, we construct photoexcited exciton wave packets by considering the internal exciton structure through first-order time-dependent perturbation theory. A zeroth-order approximation is introduced to retain the essential internal structure without evaluating the full microscopic wave function. While the internal phase contribution can be neglected for valley excitons under appropriate conditions, it becomes a relative phase between valley components for valley-mixed excitons and plays a crucial role in determining their wave-packet structures. Further numerical treatment is therefore required to fully describe valley-mixed exciton wave packets.

This work establishes a connection between structured-light angular geometry and exciton momentum-space engineering, providing a framework for controlling and probing excitonic states using tailored optical fields.
